% Bagian: sets-functions-relations
% Bab: sets
% Subbagian: schroder-bernstein

\documentclass[../../../include/open-logic-section]{subfiles}

\begin{document}

\olfileid[id]{sfr}{siz}{sb}

\olsection{Pengertian Ukuran dan Schr\"oder-Bernstein}

\begin{explain}
Berikut suatu gagasan intuitif: jika $A$ tidak lebih besar daripada $B$ dan
$B$ tidak lebih besar daripada $A$, maka $A$ dan $B$ berkardinalitas sama.
Sejujurnya, jika gagasan ini \emph{keliru}, kita hampir tidak dapat membenarkan
anggapan bahwa pengertian kesamaan kardinalitas yang kita definisikan ada
hubungannya dengan perbandingan ``ukuran'' antarkumpulan!{} Untungnya, gagasan
intuitif tersebut benar. Hal ini dibenarkan oleh Teorema Schr\"oder-Bernstein.
\end{explain}

\begin{thm}[Schr\"oder-Bernstein]
	\ollabel{thm:schroder-bernstein}
	Jika $\cardle{A}{B}$ dan $\cardle{B}{A}$,
	maka $\cardeq{A}{B}$.
\end{thm}

\begin{explain}
Dengan kata lain, jika terdapat !!a{injection} dari $A$ ke~$B$ dan
!!a{injection} dari $B$ ke~$A$, maka terdapat !!a{bijection} dari $A$ ke~$B$.

Namun, hasil ini sebenarnya cukup \emph{sulit} dibuktikan. Memang, meskipun
Cantor menyatakan hasil tersebut, orang lainlah yang membuktikannya.
\footnote{Untuk pembahasan sejarah lebih lanjut, lihat misalnya
\citet[pp.~165--6]{Potter2004}.}
\oliflabeldef{sfr:cardinals:card-sb:sec}{Kita baru akan berada dalam posisi
untuk \emph{membuktikan} Schr\"oder-Bernstein dalam
\olref[sfr][cardinals][card-sb]{sec}.}{}%
Untuk saat ini, Anda dapat (dan harus) menerimanya berdasarkan kepercayaan.

Untungnya, Schr\"oder-Bernstein \emph{benar}, dan teorema itu membenarkan cara
kita memandang relasi yang telah kita definisikan, yakni $\cardeq{A}{B}$ dan
$\cardle{A}{B}$, sebagai relasi yang berkaitan dengan ``ukuran''. Selain itu,
Schr\"oder-Bernstein sangat \emph{berguna}. Memikirkan suatu !!a{bijection}
antara dua himpunan berkardinalitas sama dapat terasa sulit. Teorema
Schr\"oder-Bernstein memungkinkan kita memecah perbandingan itu menjadi dua
arah, sehingga kita hanya perlu memikirkan !!a{injection} dari himpunan
pertama ke himpunan kedua, dan sebaliknya.
\end{explain}
\end{document}
